%% ============================================================
%%  Institutional Background
%%  Depositor Discipline in Modern Banking (2026-04-04 draft)
%%
%%  Scope: ILM/CECL transition, Day-One adjustment mechanics,
%%  treatment variable construction, adoption timeline, phase-in.
%%  Identification arguments live in sec:identification.
%% ============================================================

\section{Institutional Background}
[@@@ remove emdashes --- unless absolutely necessary, no paragraph titles]
\label{sec:background}

\subsection{From Incurred Losses to Expected Credit Losses}
\label{sec:background:transition}

U.S.\ banks historically recognized credit losses under the \textit{incurred
loss model} (ILM), codified in ASC~450-20 and ASC~310-10.  The ILM prohibited [@@@ didn't allow? check]
reserve recognition until a loss was ``probable and estimable,'' a standard
anchored to observable credit events that had already occurred --- typically a
missed payment, a borrower downgrade, or an impairment trigger --- rather than
to forward-looking assessments of latent portfolio risk.  A large body of
evidence documents the consequences: loan loss provisions lag realized defaults
by one to two years over the business cycle \citep{beatty2011delays}, and the
resulting delay in recognition reduces the informativeness of bank financial
statements precisely when outside monitors most need timely signals
\citep{bushman2012accounting, bushman2015delayed}.  From a depositor's
perspective, the ILM meant that two otherwise identical banks could carry
materially different allowances depending on portfolio vintage, borrower
geography, and the timing of internal credit reviews, none of which necessarily
reflected differences in underlying credit-risk exposure.

The 2007--2009 financial crisis concentrated attention on the procyclical
nature of delayed recognition and prompted the Financial Accounting Standards
Board (FASB) to replace the ILM trigger with a forward-looking standard.  In
June 2016, FASB issued Accounting Standards Update (ASU) 2016-13, introducing
the \textit{Current Expected Credit Loss} model (CECL), codified in ASC~326
\citep{FASB2016}.  Under CECL, a bank must estimate and reserve for
\textit{lifetime} expected credit losses at origination and at each subsequent
reporting date, drawing on historical loss experience, current conditions, and
reasonable and supportable forward-looking forecasts.  The size of the required
allowance is a direct function of the credit-risk composition of the loan
portfolio: loans with longer remaining maturities, weaker borrower quality, or
thinner collateral coverage require materially larger CECL allowances than the
ILM reserve previously carried against the same assets.

The practical significance of this shift is that CECL surfaces credit-risk
information that was economically present under the ILM but not yet required to
be disclosed.  Consistent with this, recent evidence finds that CECL adoption
increased the sensitivity of loan loss provisions to forward-looking credit
signals \citep{granja2025current, kim2026current}, reduced the procyclicality
of bank lending \citep{chen2022does}, and improved the decision-usefulness of
reserve disclosures for outside investors \citep{gee2022decision}.  Our design
exploits a specific moment in this transition --- the adoption-date catch-up
adjustment --- as a vehicle [@@@ natural experiment] for studying whether disclosed credit-risk
information disciplines uninsured depositors [@@@ and how banks respond].

\subsection{The Day-One Adjustment and Treatment Variable}
\label{sec:background:dayone}

When a bank transitions from ILM to CECL, it must reclassify its entire existing
loan portfolio under the new standard on a single date.  The difference between
the CECL-required allowance and the legacy ILM allowance --- the Day-One
adjustment --- is recorded as a cumulative-effect charge to retained earnings on
the opening balance sheet of the adoption quarter.  For bank $i$ with adoption
quarter $t^{*}$:
\begin{equation}
  \Delta_i
  \;\equiv\;
  \mathrm{ALLL}^{\mathrm{CECL}}_{i,\,t^{*}}
  -
  \mathrm{ALLL}^{\mathrm{ILM}}_{i,\,t^{*}-1},
  \label{eq:dayone}
\end{equation}
where $\mathrm{ALLL}$ denotes the allowance for loan and lease losses reported
on the Call Report.  Because CECL reserves reflect expected losses over the
full remaining contract life of each loan rather than losses already incurred,
$\Delta_i < 0$ for nearly all adopters: the new standard requires a larger
reserve than the old one.[@@@ this is not true. there ara banks with less than zero. See the first figure, and many zeros]

The Day-One charge reduces reported book equity and, absent any regulatory
accommodation, risk-weighted regulatory capital ratios by an equivalent amount.
Importantly, equation~\eqref{eq:dayone} records a pure accounting adjustment.
It does not trigger a cash transfer, dividend cut, or change in the contractual
terms of the underlying loans.  The bank's actual cash flows and its
borrowers' repayment obligations are unchanged.  What the adjustment does is
reveal, in a single mandatory disclosure, the gap between book equity as
previously reported and book equity adjusted to reflect the bank's true
forward-looking credit-loss exposure.

Our primary treatment variable scales the Day-One charge by the bank's
beginning-of-period equity:
\begin{equation}
  \texttt{CECL Adj}_{i}
  \;=\;
  \frac{\Delta_i}{\texttt{Equity}_{i,\,t-1^{*}}}
  \times 100,
  \label{eq:treatment}
\end{equation}
expressed in percentage points to facilitate comparison across banks of
different sizes.  For our main specifications we also construct a binary
indicator for banks in the top decile of the cross-sectional distribution:
\begin{equation}
  \texttt{High CECL}_{i}
  \;=\;
  \mathbf{1}\!\left\{
    \texttt{CECL Adj}_{i} \geq \tau
  \right\},
  \label{eq:binary}
\end{equation}
where $\tau$ is the top-decile threshold (~2.5 percentage points).  Results are
robust to using the continuous \texttt{CECL Adj} measure in place of the
binary indicator.

\subsection{Adoption Timeline and Regulatory Capital Phase-In}
\label{sec:background:timeline}

FASB phased the CECL transition by entity type.  Large SEC-filing registrants
--- primarily banks with assets above \$10 billion --- were required to adopt
beginning with fiscal years starting after December~15, 2019, which for
calendar-year filers means 2020Q1.  Smaller reporting companies, non-accelerated
filers, and non-public entities --- the community banks that comprise our main
sample --- adopted on January~1, 2023, and appear in their first CECL-compliant
Call Report in 2023Q1.\footnote{A small number of community banks with non-
calendar fiscal years adopted in adjacent quarters.  Our algorithm assigns the
first quarter in the 2022Q3--2023Q4 window with a non-missing
\texttt{cecl\_equity} observation as bank-specific adoption time $t^{*} = 0$.}
We focus on the 2023 community-bank cohort throughout.  The large-bank 2020
cohort serves a separate signal-validation role, discussed in
Section~\ref{sec:identification}, but its adoption coincides with the onset of
the COVID-19 pandemic, which precludes using that cohort to study deposit
market responses.

Banking regulators separately permitted community banks to elect a three-year
phase-in of the Day-One capital impact.  Under the 2020 interim final rule,
banks electing the transition could spread the
reduction to their regulatory capital ratios across 2023--2025, rather than
absorbing it fully in 2023Q1.  Two features of this election are important for
our analysis.  First, the phase-in election applied only to \textit{regulatory
capital} treatment; it had no effect on accounting recognition.  A bank
electing the phase-in still reported the full CECL reserve and the full
retained-earnings charge on its Q1 2023 Call Report, so depositors observed the
same disclosed adjustment regardless of phase-in status.  Second, the election
was made in advance of adoption, before the bank calculated the size of its
adjustment, giving it the same predetermined character as the adjustment itself.

\medskip\noindent
Three properties of the Day-One adjustment distinguish it from other credit-risk
signals and motivate our identification strategy.  The disclosure timing is set
by regulatory mandate, not bank choice; the magnitude reflects historical
portfolio composition accumulated before the standard was finalized in 2016; and
the adjustment is a pure accounting event that separates new information from
concurrent changes in cash flows or financial condition.
